\documentclass[11pt,a4paper]{article}
\usepackage[utf8]{inputenc}
\usepackage[T1]{fontenc}
\usepackage[spanish,es-minimal,es-noshorthands]{babel}
\usepackage{lmodern}
\usepackage[sfdefault,lining]{FiraSans}
\usepackage[a4paper,top=2.3cm,bottom=2.2cm,left=1.8cm,right=1.8cm,headheight=1.3cm,headsep=0.4cm,footskip=1cm]{geometry}
\usepackage{amsmath,amssymb}
\usepackage{icomma}
\usepackage{graphicx}
\usepackage[table]{xcolor}
\usepackage{tikz}
\usepackage{fontawesome5}
\usepackage{ccicons}
\usepackage{fancyhdr}
\usepackage{lastpage}
\usepackage{tasks}
\usepackage{needspace}
\usepackage[most]{tcolorbox}
\usepackage{array}
\usepackage[hidelinks]{hyperref}
\AtBeginDocument{\hypersetup{pdftitle={Ficha de ejercicios de la unidad 5: Proporcionalidad y porcentajes. Números decimales. 1º ESO},pdfsubject={Matemáticas 1º ESO},pdfkeywords={números decimales, aproximación, operaciones con decimales, decimales y fracciones, razones, proporciones, proporcionalidad directa, proporcionalidad inversa, regla de tres, porcentajes}}}

\definecolor{navy}{HTML}{1B2748}
\definecolor{accent}{HTML}{E8603A}
\definecolor{teal}{HTML}{1F9E90}
\definecolor{soft}{HTML}{EEF1F8}
\definecolor{muted}{HTML}{5D6685}

\setlength{\parindent}{0pt}
\setlength{\parskip}{3pt}
\settasks{label=\alph*), label-format=\bfseries\color{navy}, label-width=1.3em, item-indent=1.8em, column-sep=0.6em, after-item-skip=5pt, before-skip=3pt, after-skip=4pt}

% ---------- capçalera i peu ----------
\newcommand{\llicencia}{\href{https://creativecommons.org/licenses/by-nc-sa/4.0/deed.es}{\ccbyncsa}\ \ Licencia Creative Commons Atribución-NoComercial-CompartirIgual 4.0 Internacional (CC BY-NC-SA 4.0)}
\pagestyle{fancy}
\fancyhf{}
\fancyhead[L]{\raisebox{-0.35\height}{\includegraphics[height=0.95cm]{logo-gva.png}}}
\fancyhead[C]{\small\color{muted}Ficha de ejercicios · Unidad 5 · 1\textsuperscript{o} ESO}
\fancyhead[R]{\raisebox{-0.35\height}{\includegraphics[height=0.95cm]{logo-vermellar.png}}}
\fancyfoot[L]{\scriptsize\color{muted}\llicencia}
\fancyfoot[R]{\small\color{muted}\thepage}
\renewcommand{\headrulewidth}{0.4pt}
\renewcommand{\headrule}{\color{muted!40}\hrule height \headrulewidth}
\fancypagestyle{primera}{\fancyhf{}\renewcommand{\headrulewidth}{0pt}
  \fancyfoot[L]{\scriptsize\color{muted}\llicencia}
  \fancyfoot[R]{\small\color{muted}\thepage}}

% ---------- nivells i exercicis ----------
\newcommand{\estrelles}[1]{{\color{accent}\ifcase#1\or\faStar\faStar[regular]\faStar[regular]\or\faStar\faStar\faStar[regular]\or\faStar\faStar\faStar\fi}}
\newcounter{exer}
\newcommand{\ex}[1]{\par\needspace{4\baselineskip}\medskip\stepcounter{exer}\noindent{\color{navy}\bfseries\theexer.}\ \ {\small\estrelles{#1}}\ \ }

\newcommand{\seccio}[2]{\par\needspace{10\baselineskip}\bigskip
  \begin{tcolorbox}[colback=navy,colframe=navy,coltext=white,arc=2mm,boxrule=0pt,left=3mm,right=3mm,top=1.2mm,bottom=1.2mm]
  {\large\bfseries #1\quad #2}
  \end{tcolorbox}\vspace{-2pt}}

\newcolumntype{C}[1]{>{\centering\arraybackslash}p{#1}}

% quadre de recordatori (només a la fitxa d'exercicis)
\newcommand{\recorda}[1]{\par\smallskip\begin{tcolorbox}[colback=accent!7,colframe=accent!55,arc=2mm,boxrule=0.5pt,left=2.5mm,right=2.5mm,top=1mm,bottom=1mm]\small{\color{accent}\faLightbulb[regular]}\ \textbf{Recuerda:}\ #1\end{tcolorbox}}
\newcommand{\res}[1]{#1}\newcommand{\rest}[1]{#1}

% =====================================================================
% Capçalera FentMates de les fitxes: un color per a cada curs.
% Els colors són els de la pàgina inicial (LEVELS a assets/js/materials.js);
% el logo ve de _eines/logo/pdf/ (python3 _eines/logo/genera.py).
%
% Ús, al final del preàmbul de la unitat:
%   % =====================================================================
% Capçalera FentMates de les fitxes: un color per a cada curs.
% Els colors són els de la pàgina inicial (LEVELS a assets/js/materials.js);
% el logo ve de _eines/logo/pdf/ (python3 _eines/logo/genera.py).
%
% Ús, al final del preàmbul de la unitat:
%   % =====================================================================
% Capçalera FentMates de les fitxes: un color per a cada curs.
% Els colors són els de la pàgina inicial (LEVELS a assets/js/materials.js);
% el logo ve de _eines/logo/pdf/ (python3 _eines/logo/genera.py).
%
% Ús, al final del preàmbul de la unitat:
%   \input{../comu/fentmates.tex}
%   \fmcurs{1eso}
% i, al començament del document:
%   \capcalera{Matemàtiques}{1r ESO}{Fitxa d'exercicis · Unitat 2}{Divisibilitat}
% Opcional abans de \input: \def\fmneed{14} (espai mínim abans de cada secció).
% =====================================================================
\makeatletter
\def\fm@logodir{../../logo/pdf}
\providecommand{\fmneed}{10}
\def\fm@cursnom{}

% clau del curs, color i color del text damunt del color
\newcommand{\fm@curs}[3]{\@namedef{fm@color@#1}{#2}\@namedef{fm@text@#1}{#3}}
\fm@curs{1eso}{E6B84A}{navy}
\fm@curs{2eso}{3B82F6}{white}
\fm@curs{3eso}{FF7A50}{white}
\fm@curs{4eso-a}{A78BFA}{white}
\fm@curs{4eso-b}{6D28D9}{white}
\fm@curs{1batx-mi}{9B1C1C}{white}
\fm@curs{1batx-ccss}{EF6B6B}{white}
\fm@curs{2batx-mii}{1E6B47}{white}
\fm@curs{2batx-ccss}{2BB3A3}{white}

\newcommand{\fmcurs}[1]{%
  \@ifundefined{fm@color@#1}{\PackageError{fentmates}{Curs desconegut: #1}{}}{}%
  \edef\fm@tmp{\noexpand\definecolor{curs}{HTML}{\@nameuse{fm@color@#1}}%
               \noexpand\colorlet{curstext}{\@nameuse{fm@text@#1}}}\fm@tmp
  \colorlet{cursfosc}{curs!72!black}%
  \def\fm@logo{\fm@logodir/vertical-#1.pdf}%
  \def\fm@marca{\fm@logodir/marca-#1.pdf}}

% ---------- capçalera de la primera pàgina ----------
% #1 matèria (petit), #2 curs, #3 tipus de document i unitat, #4 títol
\newcommand{\capcalera}[4]{%
  \gdef\fm@cursnom{#2}%
  \par\noindent
  \begin{tcolorbox}[colback=curs,colframe=curs,boxrule=0pt,arc=2.2mm,width=3.5cm,height=2.15cm,
      valign=center,left=3mm,right=2mm,top=1mm,bottom=1mm,nobeforeafter,box align=center]
    \color{curstext}{\footnotesize #1}\par\vspace{1pt}{\fontsize{19}{21}\selectfont\bfseries #2\par}%
  \end{tcolorbox}\hspace{5mm}%
  \begin{minipage}[c]{\dimexpr\linewidth-3.5cm-5mm-2.9cm\relax}
    {\color{cursfosc}\small\bfseries\MakeUppercase{#3}}\par\vspace{2pt}
    {\color{navy}\fontsize{20}{23}\selectfont\bfseries #4\par}
  \end{minipage}\hfill
  \begin{minipage}[c]{2.6cm}\raggedleft\includegraphics[width=2.6cm]{\fm@logo}\end{minipage}\par
  \vspace{2pt}}

% ---------- barres de secció del color del curs ----------
\renewcommand{\seccio}[2]{\par\needspace{\fmneed\baselineskip}\bigskip
  \begin{tcolorbox}[colback=curs,colframe=curs,coltext=curstext,arc=2mm,boxrule=0pt,left=3mm,right=3mm,top=1.2mm,bottom=1.2mm]
  {\large\bfseries #1\quad #2}
  \end{tcolorbox}\vspace{-2pt}}

% ---------- estrelles de dificultat del color del curs (a Batxillerat no se'n posen) ----------
\@ifundefined{estrelles}{}{%
  \renewcommand{\estrelles}[1]{{\color{curs}\ifcase#1\or\faStar\faStar[regular]\faStar[regular]\or\faStar\faStar\faStar[regular]\or\faStar\faStar\faStar\fi}}}

% ---------- capçalera i peu de les pàgines ----------
% primera pàgina: logos institucionals al peu
\fancypagestyle{primera}{\fancyhf{}\renewcommand{\headrulewidth}{0pt}
  \fancyfoot[L]{\raisebox{-0.32\height}{\includegraphics[height=0.78cm]{logo-gva.png}}\hspace{5mm}%
    \raisebox{-0.32\height}{\includegraphics[height=0.78cm]{logo-vermellar.png}}}
  \fancyfoot[R]{\scriptsize\color{muted}\href{https://creativecommons.org/licenses/by-nc-sa/4.0/deed.ca}{\ccbyncsa}\ \ CC BY-NC-SA 4.0 · fentmates.org\qquad\small\thepage}}
% la resta: el símbol a l'esquerra i el curs a la dreta, amb el títol al centre com abans
\fancyhead[L]{\raisebox{-0.3\height}{\includegraphics[height=0.62cm]{\fm@marca}}}
\fancyhead[R]{\small\bfseries\color{cursfosc}\fm@cursnom}
\renewcommand{\headrule}{\color{curs!45}\hrule height \headrulewidth}
\makeatother

%   \fmcurs{1eso}
% i, al començament del document:
%   \capcalera{Matemàtiques}{1r ESO}{Fitxa d'exercicis · Unitat 2}{Divisibilitat}
% Opcional abans de \input: \def\fmneed{14} (espai mínim abans de cada secció).
% =====================================================================
\makeatletter
\def\fm@logodir{../../logo/pdf}
\providecommand{\fmneed}{10}
\def\fm@cursnom{}

% clau del curs, color i color del text damunt del color
\newcommand{\fm@curs}[3]{\@namedef{fm@color@#1}{#2}\@namedef{fm@text@#1}{#3}}
\fm@curs{1eso}{E6B84A}{navy}
\fm@curs{2eso}{3B82F6}{white}
\fm@curs{3eso}{FF7A50}{white}
\fm@curs{4eso-a}{A78BFA}{white}
\fm@curs{4eso-b}{6D28D9}{white}
\fm@curs{1batx-mi}{9B1C1C}{white}
\fm@curs{1batx-ccss}{EF6B6B}{white}
\fm@curs{2batx-mii}{1E6B47}{white}
\fm@curs{2batx-ccss}{2BB3A3}{white}

\newcommand{\fmcurs}[1]{%
  \@ifundefined{fm@color@#1}{\PackageError{fentmates}{Curs desconegut: #1}{}}{}%
  \edef\fm@tmp{\noexpand\definecolor{curs}{HTML}{\@nameuse{fm@color@#1}}%
               \noexpand\colorlet{curstext}{\@nameuse{fm@text@#1}}}\fm@tmp
  \colorlet{cursfosc}{curs!72!black}%
  \def\fm@logo{\fm@logodir/vertical-#1.pdf}%
  \def\fm@marca{\fm@logodir/marca-#1.pdf}}

% ---------- capçalera de la primera pàgina ----------
% #1 matèria (petit), #2 curs, #3 tipus de document i unitat, #4 títol
\newcommand{\capcalera}[4]{%
  \gdef\fm@cursnom{#2}%
  \par\noindent
  \begin{tcolorbox}[colback=curs,colframe=curs,boxrule=0pt,arc=2.2mm,width=3.5cm,height=2.15cm,
      valign=center,left=3mm,right=2mm,top=1mm,bottom=1mm,nobeforeafter,box align=center]
    \color{curstext}{\footnotesize #1}\par\vspace{1pt}{\fontsize{19}{21}\selectfont\bfseries #2\par}%
  \end{tcolorbox}\hspace{5mm}%
  \begin{minipage}[c]{\dimexpr\linewidth-3.5cm-5mm-2.9cm\relax}
    {\color{cursfosc}\small\bfseries\MakeUppercase{#3}}\par\vspace{2pt}
    {\color{navy}\fontsize{20}{23}\selectfont\bfseries #4\par}
  \end{minipage}\hfill
  \begin{minipage}[c]{2.6cm}\raggedleft\includegraphics[width=2.6cm]{\fm@logo}\end{minipage}\par
  \vspace{2pt}}

% ---------- barres de secció del color del curs ----------
\renewcommand{\seccio}[2]{\par\needspace{\fmneed\baselineskip}\bigskip
  \begin{tcolorbox}[colback=curs,colframe=curs,coltext=curstext,arc=2mm,boxrule=0pt,left=3mm,right=3mm,top=1.2mm,bottom=1.2mm]
  {\large\bfseries #1\quad #2}
  \end{tcolorbox}\vspace{-2pt}}

% ---------- estrelles de dificultat del color del curs (a Batxillerat no se'n posen) ----------
\@ifundefined{estrelles}{}{%
  \renewcommand{\estrelles}[1]{{\color{curs}\ifcase#1\or\faStar\faStar[regular]\faStar[regular]\or\faStar\faStar\faStar[regular]\or\faStar\faStar\faStar\fi}}}

% ---------- capçalera i peu de les pàgines ----------
% primera pàgina: logos institucionals al peu
\fancypagestyle{primera}{\fancyhf{}\renewcommand{\headrulewidth}{0pt}
  \fancyfoot[L]{\raisebox{-0.32\height}{\includegraphics[height=0.78cm]{logo-gva.png}}\hspace{5mm}%
    \raisebox{-0.32\height}{\includegraphics[height=0.78cm]{logo-vermellar.png}}}
  \fancyfoot[R]{\scriptsize\color{muted}\href{https://creativecommons.org/licenses/by-nc-sa/4.0/deed.ca}{\ccbyncsa}\ \ CC BY-NC-SA 4.0 · fentmates.org\qquad\small\thepage}}
% la resta: el símbol a l'esquerra i el curs a la dreta, amb el títol al centre com abans
\fancyhead[L]{\raisebox{-0.3\height}{\includegraphics[height=0.62cm]{\fm@marca}}}
\fancyhead[R]{\small\bfseries\color{cursfosc}\fm@cursnom}
\renewcommand{\headrule}{\color{curs!45}\hrule height \headrulewidth}
\makeatother

%   \fmcurs{1eso}
% i, al començament del document:
%   \capcalera{Matemàtiques}{1r ESO}{Fitxa d'exercicis · Unitat 2}{Divisibilitat}
% Opcional abans de \input: \def\fmneed{14} (espai mínim abans de cada secció).
% =====================================================================
\makeatletter
\def\fm@logodir{../../logo/pdf}
\providecommand{\fmneed}{10}
\def\fm@cursnom{}

% clau del curs, color i color del text damunt del color
\newcommand{\fm@curs}[3]{\@namedef{fm@color@#1}{#2}\@namedef{fm@text@#1}{#3}}
\fm@curs{1eso}{E6B84A}{navy}
\fm@curs{2eso}{3B82F6}{white}
\fm@curs{3eso}{FF7A50}{white}
\fm@curs{4eso-a}{A78BFA}{white}
\fm@curs{4eso-b}{6D28D9}{white}
\fm@curs{1batx-mi}{9B1C1C}{white}
\fm@curs{1batx-ccss}{EF6B6B}{white}
\fm@curs{2batx-mii}{1E6B47}{white}
\fm@curs{2batx-ccss}{2BB3A3}{white}

\newcommand{\fmcurs}[1]{%
  \@ifundefined{fm@color@#1}{\PackageError{fentmates}{Curs desconegut: #1}{}}{}%
  \edef\fm@tmp{\noexpand\definecolor{curs}{HTML}{\@nameuse{fm@color@#1}}%
               \noexpand\colorlet{curstext}{\@nameuse{fm@text@#1}}}\fm@tmp
  \colorlet{cursfosc}{curs!72!black}%
  \def\fm@logo{\fm@logodir/vertical-#1.pdf}%
  \def\fm@marca{\fm@logodir/marca-#1.pdf}}

% ---------- capçalera de la primera pàgina ----------
% #1 matèria (petit), #2 curs, #3 tipus de document i unitat, #4 títol
\newcommand{\capcalera}[4]{%
  \gdef\fm@cursnom{#2}%
  \par\noindent
  \begin{tcolorbox}[colback=curs,colframe=curs,boxrule=0pt,arc=2.2mm,width=3.5cm,height=2.15cm,
      valign=center,left=3mm,right=2mm,top=1mm,bottom=1mm,nobeforeafter,box align=center]
    \color{curstext}{\footnotesize #1}\par\vspace{1pt}{\fontsize{19}{21}\selectfont\bfseries #2\par}%
  \end{tcolorbox}\hspace{5mm}%
  \begin{minipage}[c]{\dimexpr\linewidth-3.5cm-5mm-2.9cm\relax}
    {\color{cursfosc}\small\bfseries\MakeUppercase{#3}}\par\vspace{2pt}
    {\color{navy}\fontsize{20}{23}\selectfont\bfseries #4\par}
  \end{minipage}\hfill
  \begin{minipage}[c]{2.6cm}\raggedleft\includegraphics[width=2.6cm]{\fm@logo}\end{minipage}\par
  \vspace{2pt}}

% ---------- barres de secció del color del curs ----------
\renewcommand{\seccio}[2]{\par\needspace{\fmneed\baselineskip}\bigskip
  \begin{tcolorbox}[colback=curs,colframe=curs,coltext=curstext,arc=2mm,boxrule=0pt,left=3mm,right=3mm,top=1.2mm,bottom=1.2mm]
  {\large\bfseries #1\quad #2}
  \end{tcolorbox}\vspace{-2pt}}

% ---------- estrelles de dificultat del color del curs (a Batxillerat no se'n posen) ----------
\@ifundefined{estrelles}{}{%
  \renewcommand{\estrelles}[1]{{\color{curs}\ifcase#1\or\faStar\faStar[regular]\faStar[regular]\or\faStar\faStar\faStar[regular]\or\faStar\faStar\faStar\fi}}}

% ---------- capçalera i peu de les pàgines ----------
% primera pàgina: logos institucionals al peu
\fancypagestyle{primera}{\fancyhf{}\renewcommand{\headrulewidth}{0pt}
  \fancyfoot[L]{\raisebox{-0.32\height}{\includegraphics[height=0.78cm]{logo-gva.png}}\hspace{5mm}%
    \raisebox{-0.32\height}{\includegraphics[height=0.78cm]{logo-vermellar.png}}}
  \fancyfoot[R]{\scriptsize\color{muted}\href{https://creativecommons.org/licenses/by-nc-sa/4.0/deed.ca}{\ccbyncsa}\ \ CC BY-NC-SA 4.0 · fentmates.org\qquad\small\thepage}}
% la resta: el símbol a l'esquerra i el curs a la dreta, amb el títol al centre com abans
\fancyhead[L]{\raisebox{-0.3\height}{\includegraphics[height=0.62cm]{\fm@marca}}}
\fancyhead[R]{\small\bfseries\color{cursfosc}\fm@cursnom}
\renewcommand{\headrule}{\color{curs!45}\hrule height \headrulewidth}
\makeatother

\fmcurs{1eso}

\begin{document}
\thispagestyle{primera}

\capcalera{Matemáticas}{1\textsuperscript{o} ESO}{Ficha de ejercicios · Unidad 5}{Proporcionalidad y porcentajes. Números decimales}
\vspace{10pt}
Nombre y apellidos: \rule{6.6cm}{0.4pt}\hfill Grupo: \rule{1.6cm}{0.4pt}\hfill Fecha: \rule{2.2cm}{0.4pt}
\vspace{8pt}

\begin{tcolorbox}[colback=soft,colframe=soft,arc=2mm,left=3mm,right=3mm,top=1.5mm,bottom=1.5mm]
\small Ejercicios graduados por niveles de dificultad:\quad \estrelles{1}\ básico \qquad \estrelles{2}\ medio \qquad \estrelles{3}\ avanzado.\\[2pt] Resuélvelos en el cuaderno mostrando todos los pasos. Las tablas y las rectas se pueden completar directamente en la ficha. Haz los cálculos sin calculadora.
\end{tcolorbox}

% =====================================================
\seccio{1}{Números decimales}

\ex{1} Escribe cómo se leen estos números:\begin{tasks}(3)
  \task $7,25$
  \task $0,408$
  \task $12,06$
  \task $3,5$
  \task $100,001$
  \task $0,09$
\end{tasks}\recorda{Leemos la parte entera, la palabra \emph{unidades} y la parte decimal con el nombre de la \textbf{última} cifra: $3,045$ son 3 unidades y 45 milésimas.}

\ex{1} Escribe con cifras:\begin{tasks}(2)
  \task Tres unidades y cuarenta y cinco milésimas
  \task Doce centésimas
  \task Veinte unidades y nueve décimas
  \task Siete milésimas
  \task Cien unidades y cinco centésimas
  \task Cuatro unidades y trescientas milésimas
\end{tasks}

\ex{1} ¿Cuánto vale la cifra indicada en cada número?\begin{tasks}(3)
  \task El 7 en $34,571$
  \task El 3 en $2,738$
  \task El 7 en $607,15$
  \task El 8 en $0,0809$
  \task El 4 en $15,4$
  \task El 6 en $9,206$
\end{tasks}

\ex{2} Completa:\begin{tasks}(2)
  \task $1$ unidad $=$ \rule{1cm}{0.4pt} décimas
  \task $1$ décima $=$ \rule{1cm}{0.4pt} centésimas
  \task $3$ unidades $=$ \rule{1cm}{0.4pt} centésimas
  \task $45$ centésimas $=$ \rule{1cm}{0.4pt} unidades
  \task $250$ milésimas $=$ \rule{1cm}{0.4pt} décimas
  \task $7$ décimas $=$ \rule{1cm}{0.4pt} milésimas
\end{tasks}

\ex{2} Descompón cada número según el valor de sus cifras:\begin{tasks}(3)
  \task $45,312$
  \task $7,09$
  \task $120,506$
\end{tasks}

\ex{3} Escribe el número que tiene: a) 4 decenas, 6 décimas y 2 milésimas; \ b) 8 centenas, 3 unidades y 5 centésimas; \ c) 15 centésimas y 3 milésimas. ¿Cuál de estos es igual a $0,15$: \ $0,150$, \ $0,015$ o \ $0,1500$?


% =====================================================
\seccio{2}{Comparar, representar y aproximar decimales}

\ex{1} Escribe $<$, $>$ o $=$:\begin{tasks}(3)
  \task $4,7\ \square\ 4,68$
  \task $0,09\ \square\ 0,1$
  \task $5,307\ \square\ 5,37$
  \task $12,5\ \square\ 12,50$
  \task $3,099\ \square\ 3,1$
  \task $0,45\ \square\ 0,405$
\end{tasks}\recorda{Para comparar, igualamos el número de decimales añadiendo ceros: $\ 3,5\ \text{?}\ 3,48\ \to\ 3,50>3,48$.}

\ex{1} Ordena de menor a mayor: \ $2,5;\ 2,05;\ 2,55;\ 2,505;\ 2,0;\ 2,15$

\ex{2} Escribe el número decimal que corresponde a cada punto:\par\smallskip{\centering \begin{tikzpicture}[baseline=0]\draw[-latex,thick] (-0.3,0)--(12.50,0);\draw[thick] (0.000,0.20)--(0.000,-0.20);\node[below=2pt,font=\small] at (0.000,-0.1) {$3$};\draw[muted] (1.200,0.12)--(1.200,-0.12);\draw[muted] (2.400,0.12)--(2.400,-0.12);\draw[muted] (3.600,0.12)--(3.600,-0.12);\draw[muted] (4.800,0.12)--(4.800,-0.12);\draw[muted] (6.000,0.12)--(6.000,-0.12);\draw[muted] (7.200,0.12)--(7.200,-0.12);\draw[muted] (8.400,0.12)--(8.400,-0.12);\draw[muted] (9.600,0.12)--(9.600,-0.12);\draw[muted] (10.800,0.12)--(10.800,-0.12);\draw[thick] (12.000,0.20)--(12.000,-0.20);\node[below=2pt,font=\small] at (12.000,-0.1) {$4$};\fill[accent] (2.400,0) circle (0.09) node[above=3pt,text=accent,font=\small]{A};\fill[accent] (8.400,0) circle (0.09) node[above=3pt,text=accent,font=\small]{B};\fill[accent] (10.800,0) circle (0.09) node[above=3pt,text=accent,font=\small]{C};\end{tikzpicture}\par}\par\smallskip{\centering \begin{tikzpicture}[baseline=0]\draw[-latex,thick] (-0.3,0)--(12.50,0);\draw[thick] (0.000,0.20)--(0.000,-0.20);\node[below=2pt,font=\small] at (0.000,-0.1) {$0,5$};\draw[muted] (1.200,0.12)--(1.200,-0.12);\draw[muted] (2.400,0.12)--(2.400,-0.12);\draw[muted] (3.600,0.12)--(3.600,-0.12);\draw[muted] (4.800,0.12)--(4.800,-0.12);\draw[muted] (6.000,0.12)--(6.000,-0.12);\draw[muted] (7.200,0.12)--(7.200,-0.12);\draw[muted] (8.400,0.12)--(8.400,-0.12);\draw[muted] (9.600,0.12)--(9.600,-0.12);\draw[muted] (10.800,0.12)--(10.800,-0.12);\draw[thick] (12.000,0.20)--(12.000,-0.20);\node[below=2pt,font=\small] at (12.000,-0.1) {$0,6$};\fill[accent] (3.600,0) circle (0.09) node[above=3pt,text=accent,font=\small]{D};\fill[accent] (9.600,0) circle (0.09) node[above=3pt,text=accent,font=\small]{E};\end{tikzpicture}\par}

\ex{2} Representa en la recta $1,4$ \ y \ $1,75$ (la recta va de 1 a 2).\par\smallskip{\centering \begin{tikzpicture}[baseline=0]\draw[-latex,thick] (-0.3,0)--(12.50,0);\draw[thick] (0.000,0.20)--(0.000,-0.20);\node[below=2pt,font=\small] at (0.000,-0.1) {$1$};\draw[muted] (0.600,0.12)--(0.600,-0.12);\draw[muted] (1.200,0.12)--(1.200,-0.12);\draw[muted] (1.800,0.12)--(1.800,-0.12);\draw[muted] (2.400,0.12)--(2.400,-0.12);\draw[muted] (3.000,0.12)--(3.000,-0.12);\draw[muted] (3.600,0.12)--(3.600,-0.12);\draw[muted] (4.200,0.12)--(4.200,-0.12);\draw[muted] (4.800,0.12)--(4.800,-0.12);\draw[muted] (5.400,0.12)--(5.400,-0.12);\draw[muted] (6.000,0.12)--(6.000,-0.12);\draw[muted] (6.600,0.12)--(6.600,-0.12);\draw[muted] (7.200,0.12)--(7.200,-0.12);\draw[muted] (7.800,0.12)--(7.800,-0.12);\draw[muted] (8.400,0.12)--(8.400,-0.12);\draw[muted] (9.000,0.12)--(9.000,-0.12);\draw[muted] (9.600,0.12)--(9.600,-0.12);\draw[muted] (10.200,0.12)--(10.200,-0.12);\draw[muted] (10.800,0.12)--(10.800,-0.12);\draw[muted] (11.400,0.12)--(11.400,-0.12);\draw[thick] (12.000,0.20)--(12.000,-0.20);\node[below=2pt,font=\small] at (12.000,-0.1) {$2$};\end{tikzpicture}\par}

\ex{2} Escribe tres números decimales comprendidos entre:\begin{tasks}(3)
  \task $0,3$ y $0,4$
  \task $2,71$ y $2,72$
  \task $5$ y $5,1$
\end{tasks}

\ex{2} Completa la tabla aproximando cada número:\par\smallskip{\centering\renewcommand{\arraystretch}{1.35}\begin{tabular}{|c|C{2.3cm}|C{2.3cm}|C{2.3cm}|}\hline Número & Truncado a las centésimas & Redondeado a las décimas & Redondeado a las centésimas\\\hline $3,4682$ & & & \\\hline $15,8962$ & & & \\\hline $0,7049$ & & & \\\hline $29,996$ & & & \\\hline\end{tabular}\par}\recorda{Para \textbf{redondear}, miramos la primera cifra que eliminamos: si es $0,1,2,3,4$ la dejamos igual; si es $5,6,7,8,9$ sumamos 1 a la última cifra que queda.}

\ex{3} Un artículo cuesta $12,4583$\,€ según la calculadora. ¿Cuánto pagaremos? Escribe un número que, redondeado a las décimas, sea $2,5$ y que sea menor que $2,5$. ¿Cuál es el número más pequeño con dos cifras decimales que redondeado a las unidades da 8?


% =====================================================
\seccio{3}{Sumar, restar y multiplicar decimales}

\ex{1} Calcula:\begin{tasks}(3)
  \task $23,45+7,8$
  \task $40-12,35$
  \task $0,9+0,35$
  \task $15-6,47$
  \task $3,06+12,9+0,4$
  \task $100-0,01$
\end{tasks}\recorda{Para sumar y restar, colocamos las \textbf{comas alineadas} y completamos con ceros.}

\ex{1} Calcula:\begin{tasks}(3)
  \task $4,6\cdot 3$
  \task $2,35\cdot 0,4$
  \task $0,6\cdot 0,05$
  \task $12,5\cdot 1,2$
  \task $3,25\cdot 1,4$
  \task $0,07\cdot 0,3$
\end{tasks}\recorda{Para multiplicar, lo hacemos como si no hubiera comas y separamos tantos decimales como \textbf{suman} los dos factores.}

\ex{1} Calcula mentalmente:\begin{tasks}(4)
  \task $3,7\cdot 100$
  \task $0,052\cdot 1000$
  \task $45,2:1000$
  \task $8:100$
  \task $6,15\cdot 10$
  \task $0,3:10$
  \task $250\cdot 0,1$
  \task $1,4\cdot 0,01$
\end{tasks}

\ex{2} Coloca la coma en el resultado sin hacer la operación: \ $3,8\cdot2,5=950$; \ $0,45\cdot1,2=540$; \ $12,6\cdot0,03=378$. Compruébalo estimando.

\ex{2} Completa con el número que falta:\begin{tasks}(3)
  \task $2,75+\square=5$
  \task $\square-1,8=3,45$
  \task $0,6\cdot\square=60$
\end{tasks}

\ex{2} Compramos $2,5$ kg de manzanas a $1,80$\,€/kg y $1,5$ kg de cerezas a $6,40$\,€/kg. Pagamos con un billete de 20\,€. ¿Cuánto nos devuelven?

\ex{3} Una piscina mide $25$ m de largo. Lucía hace $8$ largos y su hermano, $5,5$ largos. ¿Cuántos metros más ha nadado Lucía? Si Lucía nada cada largo en $0,45$ minutos, ¿cuánto tarda en total?

\ex{3} El perímetro de un triángulo es $17,25$ cm. Dos lados miden $4,8$ cm y $6,07$ cm. ¿Cuánto mide el tercer lado? ¿Cuánto mediría el perímetro de un triángulo con todos los lados 10 veces más largos?


% =====================================================
\seccio{4}{Dividir con decimales}

\ex{1} Calcula:\begin{tasks}(3)
  \task $8,4:4$
  \task $9,75:3$
  \task $6,24:3$
  \task $0,8:2$
  \task $15,3:5$
  \task $1,26:6$
\end{tasks}

\ex{1} Calcula sacando decimales en el cociente hasta que la división sea exacta:\begin{tasks}(3)
  \task $7:4$
  \task $15:8$
  \task $9:4$
  \task $3:5$
  \task $1:8$
  \task $21:12$
\end{tasks}\recorda{Si sobra resto, ponemos una coma en el cociente, añadimos un 0 al resto y seguimos dividiendo.}

\ex{2} Calcula (si el divisor tiene decimales, multiplica dividendo y divisor por 10, 100…):\begin{tasks}(3)
  \task $4,5:0,9$
  \task $2,73:0,7$
  \task $6:0,12$
  \task $12,6:0,35$
  \task $1,44:1,2$
  \task $0,9:0,03$
\end{tasks}

\ex{2} Calcula con dos cifras decimales e indica el resto: \ a) $10:3$ \quad b) $5:7$ \quad c) $25,4:6$

\ex{2} Con una botella de $1,5$ L, ¿cuántos vasos de $0,25$ L podemos llenar? ¿Y con una garrafa de 5 L?

\ex{3} Un trozo de queso de $0,85$ kg cuesta $11,39$\,€. ¿Cuál es el precio de un kilo? ¿Cuánto costaría un trozo de $1,2$ kg?

\ex{3} Un coche ha gastado $31,5$ L de gasolina en un viaje de 450 km. ¿Cuántos litros gasta cada 100 km? Si la gasolina cuesta $1,56$\,€/L, ¿cuánto ha costado el viaje?


% =====================================================
\seccio{5}{Decimales y fracciones}

\ex{1} Escribe como fracción decimal y después como fracción irreducible:\begin{tasks}(4)
  \task $0,6$
  \task $0,45$
  \task $1,2$
  \task $0,004$
  \task $0,125$
  \task $2,75$
  \task $0,08$
  \task $3,5$
\end{tasks}\recorda{En el numerador ponemos el número sin coma; en el denominador, un 1 seguido de tantos ceros como decimales tiene: \ $0,35=\dfrac{35}{100}=\dfrac{7}{20}$.}

\ex{1} Escribe en forma decimal:\begin{tasks}(4)
  \task $\dfrac{7}{10}$
  \task $\dfrac{35}{100}$
  \task $\dfrac{4}{1000}$
  \task $\dfrac{125}{100}$
\end{tasks}

\ex{2} Escribe en forma decimal y clasifica en exacto, periódico puro o periódico mixto:\begin{tasks}(4)
  \task $\dfrac{9}{20}$
  \task $\dfrac{5}{9}$
  \task $\dfrac{7}{12}$
  \task $\dfrac{3}{8}$
  \task $\dfrac{4}{11}$
  \task $\dfrac{1}{6}$
  \task $\dfrac{13}{5}$
  \task $\dfrac{2}{3}$
\end{tasks}\recorda{Dividimos el numerador entre el denominador. Si las cifras se repiten, el decimal es \textbf{periódico} y marcamos el período: \ $\dfrac{5}{6}=0,8333\ldots=0,8\overline{3}$.}

\ex{2} Ordena de menor a mayor pasándolo todo a decimal: \ $0,7;\ \ \dfrac{2}{3};\ \ 0,67;\ \ \dfrac{3}{4};\ \ 0,6$

\ex{2} ¿Cuál es mayor?\begin{tasks}(3)
  \task $\dfrac{5}{8}$ o $0,62$
  \task $\dfrac{1}{3}$ o $0,33$
  \task $\dfrac{7}{5}$ o $1,4$
\end{tasks}

\ex{3} Escribe una fracción de denominador 3 comprendida entre $1,5$ y $2$, y una fracción de denominador 8 comprendida entre $0,3$ y $0,4$.

\ex{3} Un pastel se divide en 8 porciones iguales. Marta se come $0,25$ del pastel y Pablo 3 porciones. ¿Quién ha comido más? ¿Qué parte del pastel queda, en fracción y en decimal?


% =====================================================
\seccio{6}{Razones y proporciones}

\ex{1} En una cesta hay 18 manzanas y 24 peras. Calcula la razón entre manzanas y peras y la razón entre peras y frutas totales. ¿Qué significan?

\ex{1} Comprueba si forman proporción:\begin{tasks}(4)
  \task $\dfrac{6}{9}$ y $\dfrac{8}{12}$
  \task $\dfrac{4}{10}$ y $\dfrac{6}{15}$
  \task $\dfrac{2,5}{4}$ y $\dfrac{5}{9}$
  \task $\dfrac{3}{7}$ y $\dfrac{12}{28}$
\end{tasks}\recorda{$\dfrac{a}{b}=\dfrac{c}{d}$ forman proporción si $a\cdot d=b\cdot c$.}

\ex{2} Calcula el término desconocido:\begin{tasks}(3)
  \task $\dfrac{3}{5}=\dfrac{12}{x}$
  \task $\dfrac{x}{8}=\dfrac{15}{24}$
  \task $\dfrac{4}{x}=\dfrac{10}{25}$
  \task $\dfrac{2,5}{4}=\dfrac{x}{12}$
  \task $\dfrac{1,2}{0,3}=\dfrac{8}{x}$
  \task $\dfrac{7}{x}=\dfrac{21}{6}$
\end{tasks}

\ex{2} Para hacer una masa de pizza hacen falta 300 g de harina por cada 180 mL de agua. ¿Cuánta agua necesitamos para 750 g de harina?

\ex{3} La razón entre la edad de una madre y la de su hija es $\dfrac{7}{2}$. Si la hija tiene 12 años, ¿cuántos tiene la madre? ¿Y cuántos años tenía la madre cuando la hija tenía 4? ¿Se mantiene la razón?


% =====================================================
\seccio{7}{Proporcionalidad directa}

\ex{1} Completa la tabla de proporcionalidad directa y calcula la constante de proporcionalidad:\par\smallskip{\centering {\renewcommand{\arraystretch}{1.35}\begin{tabular}{|l|C{1.25cm}|C{1.25cm}|C{1.25cm}|C{1.25cm}|C{1.25cm}|}\hline Entradas & 1 & 2 & 4 &  & 10\\\hline Precio (€) &  & 13 &  & 45,50 & \\\hline\end{tabular}}\par}

\ex{1} Indica si son directamente proporcionales:\begin{tasks}(2)
  \task Las barras de pan que compramos y el precio que pagamos.
  \task La edad de una persona y su altura.
  \task El lado de un cuadrado y su perímetro.
  \task Los litros de gasolina y los km recorridos (con el mismo coche).
\end{tasks}

\ex{1} Si 4 bolígrafos cuestan $3,20$\,€, ¿cuánto cuestan 9 bolígrafos? ¿Y 15?

\ex{2} Un coche gasta $6,5$ L de gasolina cada 100 km. ¿Cuántos litros gastará en un viaje de 340 km?

\ex{2} Un grifo llena 120 L en 8 minutos. ¿Cuánto tardará en llenar un depósito de 450 L? ¿Cuántos litros habrá echado en media hora?

\ex{2} Una receta de bizcocho para 6 personas lleva 250 g de harina, 4 huevos y $0,15$ L de leche. ¿Qué cantidades necesitamos para 9 personas?

\ex{3} En un mapa, 2 cm representan 15 km reales. ¿Qué distancia real hay entre dos ciudades que en el mapa están a $9,4$ cm? ¿A cuántos cm del mapa están dos pueblos separados por 6 km?

\ex{3} Un taxi cobra 3\,€ por subir y $1,20$\,€ por km. Completa la tabla y di si el precio es directamente proporcional a los km.\par\smallskip{\centering {\renewcommand{\arraystretch}{1.35}\begin{tabular}{|l|C{1.25cm}|C{1.25cm}|C{1.25cm}|C{1.25cm}|}\hline km & 1 & 2 & 5 & 10\\\hline Precio (€) &  &  &  & \\\hline\end{tabular}}\par}


% =====================================================
\seccio{8}{Proporcionalidad inversa}

\ex{1} Completa la tabla de proporcionalidad inversa y calcula la constante:\par\smallskip{\centering {\renewcommand{\arraystretch}{1.35}\begin{tabular}{|l|C{1.25cm}|C{1.25cm}|C{1.25cm}|C{1.25cm}|C{1.25cm}|}\hline Velocidad (km/h) & 40 & 60 & 80 &  & 120\\\hline Tiempo (h) & 6 &  &  & 2,5 & \\\hline\end{tabular}}\par}

\ex{1} Indica si son directamente proporcionales, inversamente proporcionales o ninguna de las dos cosas:\begin{tasks}(2)
  \task El número de obreros y el tiempo que tardan en hacer una pared.
  \task Los litros de pintura y los m$^2$ pintados.
  \task El número de amigos que pagan un regalo y lo que paga cada uno.
  \task El peso de una persona y su edad.
\end{tasks}

\ex{1} Si 3 grifos llenan una piscina en 10 horas, ¿cuánto tardarán 5 grifos iguales? ¿Y 2 grifos?

\ex{2} Un grupo de 12 personas tiene comida para 30 días. Si se añaden 8 más, ¿para cuántos días tendrán comida?

\ex{2} Un tren a 90 km/h tarda 2 horas en hacer un trayecto. ¿Cuánto tardaría a 120 km/h? ¿A qué velocidad tendría que ir para tardar 3?

\ex{3} 5 máquinas hacen un trabajo en 12 horas. ¿Cuántas máquinas hacen falta para hacerlo en 4 horas? Si solo tenemos 3, ¿cuánto tardarán?

\ex{3} Para repartir el agua de un depósito hacen falta 60 garrafas de 5 L. ¿Cuántas garrafas de $1,5$ L harían falta? Si solo tenemos 25 garrafas iguales, ¿de qué capacidad tendrían que ser?


% =====================================================
\seccio{9}{Porcentajes}

\ex{1} Completa la tabla:\par\smallskip{\centering\renewcommand{\arraystretch}{1.5}\begin{tabular}{|c|C{2.4cm}|C{2.4cm}|}\hline Porcentaje & Fracción irreducible & Decimal\\\hline $40\,\%$ & & \\\hline $5\,\%$ & & \\\hline $75\,\%$ & & \\\hline $12\,\%$ & & \\\hline $120\,\%$ & & \\\hline $2,5\,\%$ & & \\\hline\end{tabular}\par}

\ex{1} Calcula:\begin{tasks}(3)
  \task El $20\,\%$ de $45$
  \task El $75\,\%$ de $120$
  \task El $15\,\%$ de $60$
  \task El $8\,\%$ de $350$
  \task El $12\,\%$ de $250$
  \task El $7,5\,\%$ de $60$
\end{tasks}\recorda{El $p\,\%$ de $N$ es $N\cdot\dfrac{p}{100}$. Por ejemplo, el $35\,\%$ de $80$ es $80\cdot0,35=28$.}

\ex{1} Expresa en porcentaje:\begin{tasks}(6)
  \task $0,35$
  \task $0,06$
  \task $1,5$
  \task $\dfrac{3}{5}$
  \task $\dfrac{7}{20}$
  \task $\dfrac{1}{8}$
\end{tasks}

\ex{2} En un instituto de 650 alumnos, 91 vienen en bicicleta. ¿Qué porcentaje representan? ¿Y en una clase de 25 alumnos, si 6 llevan gafas?

\ex{2} En una clase de 28 alumnos, el $25\,\%$ ha sacado sobresaliente y el $50\,\%$ notable. ¿Cuántos alumnos han sacado cada nota? ¿Qué porcentaje no ha sacado ni sobresaliente ni notable?

\ex{3} Hemos leído 135 páginas de un libro, que son el $45\,\%$ del total. ¿Cuántas páginas tiene el libro?

\ex{3} ¿Qué es más, el $30\,\%$ de 50 o el $50\,\%$ de 30? ¿Y el $20\,\%$ de 80 o el $80\,\%$ de 20? Explica por qué pasa.


% =====================================================
\seccio{10}{Aumentos y disminuciones porcentuales}

\ex{1} Escribe el índice de variación:\begin{tasks}(3)
  \task aumento del 30 \%
  \task disminución del 8 \%
  \task aumento del 4 \%
  \task rebaja del 50 \%
  \task IVA del 21 \%
  \task descuento del 15 \%
\end{tasks}\recorda{Aumento del $p\,\%$: multiplicamos por $1+\dfrac{p}{100}$. \ Disminución del $p\,\%$: multiplicamos por $1-\dfrac{p}{100}$.}

\ex{1} Calcula la cantidad final:\begin{tasks}(2)
  \task $80$\,€ con un descuento del $15\,\%$
  \task $50$\,€ con una subida del $21\,\%$
  \task $650$\,€ con un descuento del $12\,\%$
  \task $18$\,€ con una subida del $4\,\%$
  \task $2400$\,€ con una subida del $5\,\%$
  \task $64,90$\,€ con un descuento del $30\,\%$
\end{tasks}

\ex{2} Una camiseta pasa de 25\,€ a 20\,€. ¿Qué porcentaje de descuento le han hecho? ¿Y si un abrigo pasa de 120\,€ a 138\,€, qué porcentaje ha subido?

\ex{2} En una tienda, un televisor cuesta 450\,€ sin IVA. ¿Cuánto cuesta con el IVA del 21\,\%? Si además nos hacen un descuento del 10\,\% sobre el precio con IVA, ¿cuánto pagamos?

\ex{3} Después de aplicar el IVA del 21\,\%, una factura es de $108,90$\,€. ¿Cuánto valía sin IVA? ¿Y si después de una rebaja del 15\,\% pagamos 68\,€, cuál era el precio inicial?

\ex{3} El precio de un patinete de 200\,€ sube un 10\,\% y, al cabo de un mes, baja un 10\,\%. ¿Cuánto cuesta al final? ¿Vuelve al precio inicial? ¿Qué pasaría si primero bajara un 20\,\% y después subiera un 20\,\%?

\ex{3} La población de un pueblo era de 2400 habitantes. El primer año aumenta un 5\,\% y el segundo año disminuye un 5\,\%. ¿Cuántos habitantes tiene al final? ¿Qué porcentaje ha variado respecto al inicio?

\vfill
\begin{center}\color{muted}\small\faCheckCircle[regular]\ \ ¡Ánimo! Si te cuesta algún ejercicio, repasa los apuntes interactivos de la unidad.\end{center}

\end{document}
